% capitulos/02-materiais-metodos.tex

Lorem ipsum dolor sit amet, consectetur adipiscing elit, sed do eiusmod
tempor incididunt ut labore et dolore magna aliqua \cite{exemplo:artigo:2024}.
Os procedimentos adotados nesta pesquisa seguem as diretrizes estabelecidas
por \citeonline{exemplo:livro:2020}.

\section{Área de Estudo}

A área de estudo compreende uma parcela de 50 ha inserida na Estação
Ecológica de Águas Emendadas (\abreviatura{ESECAE}{Estação Ecológica de
Águas Emendadas}), localizada no Distrito Federal (DF), Brasil
(15°32'S, 47°36'W, altitude média de 1.050 m). O clima da região é
classificado como Aw (tropical savânico) segundo a classificação de
Köppen-Geiger, com precipitação média anual de 1.500 mm e temperatura
média de 22 °C \cite{exemplo:web:2024}.

\begin{figure}[H]
  \centering
  \fbox{\rule{0pt}{6cm}\rule{12cm}{0pt}}
  \caption{Pontos de coleta de amostras de solo e serapilheira na
           área de estudo (ESECAE, DF). Os círculos indicam parcelas
           amostradas e os triângulos indicam pontos controle.
           Fonte: elaborado pelo autor (2025).}
  \label{fig:pontos-coleta}
\end{figure}

\section{Coleta e Processamento de Amostras}

As coletas foram realizadas em três campanhas sazonais: período seco
(julho/2024), transição (outubro/2024) e período chuvoso (janeiro/2025).
Em cada campanha foram amostrados 30 pontos distribuídos sistematicamente
conforme a \autoref{fig:pontos-coleta}. Em cada ponto foram coletadas
amostras de solo nas profundidades de 0--10 cm, 10--20 cm e 20--30 cm,
totalizando 270 amostras de solo por campanha.

As amostras foram acondicionadas em sacos plásticos de polietileno,
identificadas e transportadas sob refrigeração ($4\,^\circ\text{C}$) ao
\abreviatura{LQAA}{Laboratório de Química Ambiental e Agroecologia}
da Faculdade UnB de Planaltina.

\section{Análises Laboratoriais}

\subsection{Determinação de Mercúrio Total}

O mercúrio total (\abreviatura{HgT}{Mercúrio Total}) foi determinado
por espectrometria de absorção atômica com vapor frio (\abreviatura{CV-AAS}{%
Cold Vapor Atomic Absorption Spectrometry}), utilizando espectrômetro
modelo RA-915+ (Lumex, Rússia), conforme metodologia descrita por
\citeonline{exemplo:artigo:2024}. O limite de detecção do método foi
de 0,001 ng g$^{-1}$ (peso seco).

\subsection{Análises Físico-Químicas do Solo}

O pH foi determinado em solução aquosa (relação solo:água = 1:2,5).
O teor de carbono orgânico total (\abreviatura{COT}{Carbono Orgânico Total})
foi determinado por oxidação via úmida com dicromato de potássio em meio
ácido \cite{exemplo:livro:2020}. A \autoref{tab:analises} apresenta
os parâmetros analisados e os respectivos métodos.

\begin{table}[H]
  \centering
  \caption{Parâmetros físico-químicos analisados nas amostras de solo
           e respectivos métodos de análise.}
  \label{tab:analises}
  \begin{tabularx}{\textwidth}{lXl}
    \toprule
    \textbf{Parâmetro} & \textbf{Método} & \textbf{Referência} \\
    \midrule
    pH                 & Potenciometria (solo:água 1:2,5)       & \cite{exemplo:livro:2020} \\
    \thinrule
    COT (\%)           & Oxidação via úmida (Walkley-Black)     & \cite{exemplo:livro:2020} \\
    \thinrule
    HgT (ng g$^{-1}$) & CV-AAS (RA-915+)                       & \cite{exemplo:artigo:2024} \\
    \thinrule
    Granulometria      & Pipeta (EMBRAPA)                       & \cite{exemplo:relatorio:2021} \\
    \thinrule
    Umidade (\%)       & Gravimetria ($105\,^\circ\text{C}$, 24 h)          & \cite{exemplo:livro:2020} \\
    \bottomrule
  \end{tabularx}
  \fonte{Elaborado pelo autor (2025).}
\end{table}

\section{Análise Estatística}

Os dados foram submetidos ao teste de normalidade de Shapiro-Wilk
($\alpha = 0{,}05$). Para comparação entre campanhas e profundidades,
utilizou-se análise de variância (\abreviatura{ANOVA}{Analysis of
Variance}) de dois fatores, seguida do teste post-hoc de Tukey
($p < 0{,}05$). O fator de bioacumulação (\abreviatura{BAF}{Bioaccumulation
Factor}) foi calculado pela \autoref{eq:baf}:

\begin{equation}
  \label{eq:baf}
  \mathrm{BAF} = \frac{C_{\text{organismo}}}{C_{\text{solo}}}
\end{equation}

\noindent em que $C_{\text{organismo}}$ é a concentração de HgT no
organismo (ng g$^{-1}$, peso seco) e $C_{\text{solo}}$ é a
concentração de HgT no solo (ng g$^{-1}$, peso seco).

Todas as análises estatísticas foram realizadas no ambiente R versão
4.3.2 \cite{exemplo:web:2024}, com nível de significância de 5\%.

